\documentclass[11pt,a4paper]{article}  % Déclare la classe du document.

\newcommand{\chaptermark}[2]{\og #2 \fg{} (#1)} % ligne uniquement pour éviter une erreur en cas d'article
\include{../1ere-preambule}

\usepackage{epic}   % extension et simplification de  l'environnement picture
\usepackage{graphicx}    % permet d'inclure des graphique externe dans un document
\usepackage{arydshln}   %permet d'avoir des lignes en pointillés

\newcommand{\lignes}[1]{
	\setlength{\unitlength}{1mm}
	\vskip 0,4cm\noindent
	\multido{\i=0+1}{#1}{%
		\noindent
		\begin{picture}(150,5)
			\linethickness{0,005mm}
			\put(-5,0){\dashline{1}(175,5)(5,5)}
		\end{picture}\vskip 0,001cm
	}
	\vskip 0cm
}

\rhead{DS n$^{\circ}02$ - Probabilités - 02/05/2023}

\newcommand{\va}[1]{\lvert \ #1 \ \rvert}
\newcommand{\vaf}[1]{\left\lvert \ #1\  \right\rvert}

\begin{document}
	\graphicspath{{images/}}
	
	\begin{center}
		
		
		\fbox{
			\begin{minipage}[]{15cm}
				\begin{center}
					\LARGE{\rule{0.5cm}{0pt}\rule[-0.25cm]{0pt}{0.9cm}
						DS n$^{\circ}02$ \\
						Probabilités 
					}
				\end{center}
			\end{minipage}
		}
	\end{center}

\begin{center}1h20min - Barème indicatif\\
%\textbf{Les élèves avec un tier-temps ne traitent pas les questions avec  le symbole \ding{46}}\\
Calculatrice autorisée.\\Les résultats devront être justifiés (à l'aide de calcul ou de propriétés)
\end{center}


\exerciceDS{ - (4 points)}
%environ 15-20 min
Le tableau ci-dessous donne les résultats d'une étude des plats choisis par 200 élèves à la cantine d'un lycée.\\
\begin{tabular}{|l|Xc|Xc|Xc|}
	\cline { 2 - 4 } \multicolumn{1}{Xc|}{} & Frites & Haricots & TOTAL \\
	\hline Compote & 90 & 70 & 160 \\
	\hline Crème Dessert & 30 & 10 & 40 \\
	\hline Total & 120 & 80 & 200 \\
	\hline
\end{tabular}
\begin{enumMet}
	\item Calculer la fréquence marginale des deux sous-populations suivantes:
	\begin{enumMet}
		\item élèves ayant pris des frites;
		\item élèves ayant pris une crème dessert.
	\end{enumMet}
	\item Calculer la fréquence conditionnelle des deux sous-populations suivantes:
	\begin{enumMet}
		\item élèves ayant pris une compote parmi ceux qui ont pris des frites:
		\item élèves ayant pris des haricots parmi ceux qui ont pris une crème dessert.
	\end{enumMet}
\end{enumMet}


\exerciceDS{ - (6 points)}
Une imprimerie produit des journaux sur deux machines d'impression A et B.\\
On relève, sur un échantillon de 800 journaux, le nombre de journaux présentant au moins un défaut.\\
\begin{tabular}{|l|Xc|Xc|Xc|}
	\cline { 2 - 4 } \multicolumn{1}{Xc|}{} & Machine A & Machine B & Total \\
	\hline Sans défaut & 494 & 266 & 760 \\
	\hline Avec défaut & 26 & 14 & 40 \\
	\hline Total & 520 & 280 & 800 \\
	\hline
\end{tabular}


On prélève au hasard un journal de cet échantillon. On note A l'évènement \guillemets{le journal a été imprimé sur la machine $A$} et $D$ l'évènement \guillemets{le journal présente un défaut}.
\begin{enumMet}
	\item Calculer $P(A), P(A \cap D)$ et $P_A(D)$.
	\item Recopier et compléter les deux arbres de probabilités ci-dessous.
	\begin{multicols}{2}
		\Proba[Arbre]{$A$/...,$\overline{A}$/...,$D$/...,$\overline{D}$/...,$D$/...,$\overline{D}$/...}
		
		\Proba[Arbre]{$D$/...,$\overline{D}$/...,$A$/...,$\overline{A}$/...,$A$/...,$\overline{A}$/...}
	\end{multicols}
	\item Les évènements A et D sont-ils indépendants ? Interpréter la réponse.
\end{enumMet}

\exerciceDS{ - (5 points)}
Une expérience aléatoire admet deux issues que l'on note $S$ \guillemets{succès} et $E$ \guillemets{échec}. On suppose que $P(S)=0,2$. On réalise successivement deux fois la même expérience aléatoire et de manière indépendante.
\begin{enumMet}
	\item Quelle est la probabilité que la seconde expérience soit un succès sachant que la première est un succès?
	\item Recopier et compléter l'arbre de probabilités ci-dessous.\\
	\Proba[Arbre]{$S$/...,$E$/...,$S$/...,$E$/...,$S$/...,$E$/...}
	\item Calculer la probabilité d'obtenir deux succès successifs.
	\item Calculer la probabilité que la première expérience soit un échec et que la seconde soit un succès.
	\item Calculer la probabilité d'obtenir un seul succès sur les deux expériences.
\end{enumMet}


\exerciceDS{ - (3 points)}
Les réservations d'un hôtel se font soit par son site Internet, soit auprès d'une agence.\\
Une étude réalisée par l'hôtel montre que pour une réservation auprès d'une agence, $5 \%$ des clients ne se présentent pas à l'hôtel, alors que pour une réservation sur son site, $2 \%$ ne se présentent pas. $30 \%$ des réservations sont effectuées auprès d'une agence.\\
On choisit une réservation au hasard et on considère les évènements $A$ \guillemets{la réservation a été faite en agence} et $\mathrm{H}$ \guillemets{le client se présente à l'hôtel}.\\
\begin{enumMet}
	\item Recopier et compléter l'arbre de probabilités ci-dessous.\\
	\Proba[Arbre]{$A$/...,$\overline{A}$/...,$H$/...,$\overline{H}$/...,$H$/...,$\overline{H}$/...}
	\item Calculer la probabilité qu'un client ne se présente pas à l'hôtel (Calculer $P(\overline{H})$).
\end{enumMet}
D'après le Bac STHR, Métropole, 2020.


\exerciceDS{ - (3 points)}
A et $B$ sont deux évènements tels que $P(A)=0,4$; $P(B)=0,9$ et $P(A \cap B)=0,36$.
\begin{enumMet}
	\item Calculer $P_A(B)$ et $P_B(A)$.
	\item Les évènements $A$ et $B$ sont-ils indépendants ?
\end{enumMet}
\end{document}

